\section{Desafios e oportunidades atuais}

\subsection{Evitando colisões}

O uso de algoritmos de desvio e planejamento de trajetória dinâmica é um elemento importante para garantir a segurança das missões espaciais, especialmente em ambientes complexos e dinâmicos, como as órbitas terrestres, operações entre espaçonaves ou próximas a corpos celestes.
Esses algoritmos permitem que espaçonaves comuniquem entre si, detectem e evitem potenciais colisões com detritos espaciais, outros veículos espaciais ou estruturas orbitais; sendo assim, elas ajustam suas trajetórias de forma autônoma, sem depender de comandos diretos de humanos, como explicado por \cite{newman2024}.
O desenvolvimento e a integração desses algoritmos suportarão as futuras missões ao redor da Lua e de Marte.

\subsection{Autonomia total ou Automação avançada}

Os sistemas com automação avançada e sistemas totalmente autônomos possuem diferenças no quesito de dependência de intervenção humana e no grau de controle e decisão necessária para cada situação encontrada.

\subsubsection{Sistemas com automação avançada}

Como explicado por \cite{ibmAutomation}, os sistemas com automação avançada são sistemas projetados para executar tarefas complexas com a mínima intervenção humana, no entanto, dependem de supervisão ou autorização para certas decisões críticas.
No contexto espacial, os sistemas com automação avançada podem executar etapas de navegação, controle e monitoramento de forma independente, mas requerem validação ou comando final de operadores humanos \cite{nasaNAC2018}, principalmente em situações de risco ou de incerteza.

\subsubsection{Sistemas totalmente autônomos}

Os sistemas totalmente autônomos, conforme explica \cite{ibmAutomation}, operam sem qualquer intervenção humana, desde a análise do ambiente até a tomada de decisões e execução de ações.
Eles são projetados para avaliar situações, identificar riscos e ajustar suas operações em tempo real, mesmo em cenários inesperados.
No contexto espacial, isso significa que uma espaçonave totalmente autônoma pode planejar e executar manobras, evitar colisões e realizar acoplamentos sem a necessidade de comunicação com operadores terrestres ou astronautas a bordo.
Essa autonomia completa é útil em missões em que a latência da comunicação torna o controle em tempo real impraticável, como em operações ao redor de Marte ou em exploração em espaço profundo.

\subsection{Barreiras Tecnológicas}

Atualmente, a transição de sistemas automatizados para sistemas totalmente autônomos enfrenta desafios tecnológicos.
Desenvolver algoritmos que lidam com a complexidade e imprevisibilidade de ambientes dinâmicos no espaço exigem sensores avançados integrados com processamento de dados em tempo real, e que tem sido revelado como custoso computacionalmente falando \cite{ainowComputeAI}.
Esses sistemas devem ser robustos o suficiente para operar de forma confiável em condições adversas, como as temperaturas extremas, radiação e a ausência de gravidade.
A maior preocupação é a confiabilidade, pois falhas nesse tipo de sistema podem comprometer missões inteiras sem possibilidade de qualquer intervenção humana.

\subsection{Barreiras Regulatórias}

Além dos desafios tecnológicos, a transição para sistemas autônomos enfrenta barreiras regulatórias, como a legislação espacial atual foi desenvolvida com base em operações controladas por humanos, além dos tratados internacionais, como o Tratado do Espaço de 1967, não contemplam explicitamente a responsabilidade e os critérios de segurança para sistemas que operam sem supervisão humana.
As questões sobre a responsabilidade em caso de falhas, ética no uso de autonomia \cite{ainowComputeAI} e a interoperabilidade entre diferentes sistemas autônomos ainda precisam ser discutidas e definidas.
Além disso, a falta de padrões técnicos globais dificulta a integração e o uso seguro desses sistemas em missões multinacionais.

\subsection{Interoperabilidade global}

O aumento da presença de empresas privadas e de novos países no setor espacial tem transformado o panorama global da exploração e utilização do espaço.
Empresas privadas, como a \emph{SpaceX}, \emph{Blue Origin} e \emph{Rocket Lab}, têm expandido suas operações oferecendo serviços de lançamento, desenvolvimento de satélites e missões tripuladas e não tripuladas.
Paralelamente, países como Emirados Árabes Unidos, Índia e Brasil \cite{senadoALADA} estão aumentando seus investimentos em tecnologias espaciais e estabelecendo metas para explorar o espaço.
Essa diversificação é benéfica, pois acelera a inovação e reduz os custos de acesso ao espaço \cite{jones2018}.
Contudo, ela também introduziu novos desafios, principalmente no que se refere à padronização de tecnologias e protocolos.

A criação de padrões universais para interfaces \cite{nasaManualOfDocking2017} de acoplamento é uma das prioridades no cenário espacial atual.
Com a multiplicação de novos operadores, a falta de padronização pode resultar em incompatibilidades que dificultam operações conjuntas, aumentam os custos e comprometem a segurança.
Essa introdução de padrões universais oferece benefícios, como:

\begin{enumerate}[label=(\roman*)]
    \item Interoperabilidade: Permite que espaçonaves de diferentes países ou de empresas privadas possam se conectar sem dificuldades.
    \item Redução de custos: Evita a necessidade de desenvolvimento contínuo de tecnologias exclusivas ou adaptadoras.
    \item Segurança e confiabilidade: Permite que os sistemas sejam compatíveis e testados em cenários diversos, permitindo a minimização dos riscos.
    \item Incentivo à cooperação internacional: Promove um ambiente de colaboração no espaço para resolver desafios globais, como a exploração de Marte ou missões em espaço profundo.
\end{enumerate}